% reports/ieee_report/sections/p1_baselines.tex
% Author: Vani Rudra (Task 03 — Classical ML Baselines)
% Phase 1: Machine Learning only.
%
% Describes the classical anomaly-detection baselines used as fair
% comparisons against the GMM density estimator.

\subsection{Classical Anomaly-Detection Baselines}
\label{sec:baselines}

To provide a rigorous fair comparison against the Gaussian Mixture Model
(Sec.~\ref{sec:gmm}), we implement four classical anomaly detectors plus
two reference baselines.  All models are trained exclusively on
\textbf{legitimate training transactions}; fraud labels are used only for
hyper-parameter selection on the validation split and for final evaluation
on the held-out test split.

% ─── Isolation Forest ────────────────────────────────────────────────────
\subsubsection{Isolation Forest}
\label{sec:iforest}

The Isolation Forest~\cite{liu2008isolation} exploits the observation that
anomalies are ``few and different'': a random recursive partitioning tree
isolates them in fewer splits.  For each sample~$\mathbf{x}$, define
$h(\mathbf{x})$ as the path length (number of splits) required to isolate
it.  The anomaly score is
\begin{equation}
    s(\mathbf{x}, n) = 2^{-E[h(\mathbf{x})]\,/\,c(n)},
    \label{eq:iforest_score}
\end{equation}
where $c(n) = 2H(n-1) - 2(n-1)/n$ is the average path length of an
unsuccessful search in a BST of $n$ elements and $H(\cdot)$ is the harmonic
number.  Scores close to~1 indicate anomalies; scores close to~0.5 indicate
normal points.

We search over \texttt{n\_estimators} $\in \{100, 200, 500\}$,
\texttt{max\_samples} $\in \{256, 1024, \text{auto}\}$, and
\texttt{max\_features} $\in \{0.5, 1.0\}$, selecting the configuration that
maximises validation PR-AUC.  Contamination is \emph{not} used during fitting.

% ─── PCA Reconstruction Error ────────────────────────────────────────────
\subsubsection{PCA Reconstruction Error}
\label{sec:pca_recon}

Given the top-$k$ principal components $\mathbf{U}_k$ of the legitimate
training covariance, the projection $\hat{\mathbf{x}} =
\mathbf{U}_k \mathbf{U}_k^\top \mathbf{x}$ reconstructs the part of
$\mathbf{x}$ in the principal subspace.  The anomaly score is the squared
reconstruction error:
\begin{equation}
    s(\mathbf{x}) = \|\mathbf{x} - \hat{\mathbf{x}}\|^2
                  = \mathbf{x}^\top (\mathbf{I} - \mathbf{U}_k \mathbf{U}_k^\top) \mathbf{x}.
    \label{eq:pca_recon}
\end{equation}
This equals the squared Mahalanobis distance in the \emph{minor}
(residual) subspace when the covariance eigenvalues are uniform there, thus
connecting PCA reconstruction to probabilistic outlier detection.

We sweep $k$ from~1 to~31 (the full feature dimensionality), also noting
the $k$ values required for 90\%, 95\% and 99\% explained variance.  The
value of~$k$ that maximises validation PR-AUC is selected.

% ─── One-Class SVM ───────────────────────────────────────────────────────
\subsubsection{One-Class SVM (RBF Kernel)}
\label{sec:ocsvm}

The One-Class SVM~\cite{scholkopf2001estimating} maps data into a
high-dimensional feature space via the RBF kernel
$K(\mathbf{x},\mathbf{x}') = \exp(-\gamma\|\mathbf{x}-\mathbf{x}'\|^2)$,
then finds the maximum-margin hyperplane separating the data from the origin.
The parameter $\nu \in (0,1]$ provides an upper bound on the fraction of
outliers and a lower bound on the fraction of support vectors.

The anomaly score is $s(\mathbf{x}) = -f(\mathbf{x})$, where $f$ is the
signed distance to the decision boundary (positive for inliers).  We search
$\nu \in \{0.001, 0.01, 0.05\}$ and $\gamma \in \{\text{scale}, 0.01, 0.1\}$.

Because OC-SVM has $O(n^2 d)$ to $O(n^3 d)$ fitting complexity, we fit on a
seeded random sub-sample of 40{,}000 legitimate training rows.  This is a
standard trade-off between computational cost and model fidelity.

% ─── Local Outlier Factor ────────────────────────────────────────────────
\subsubsection{Local Outlier Factor}
\label{sec:lof}

LOF~\cite{breunig2000lof} computes a local reachability density for each point
relative to its $k$-nearest neighbours.  The LOF score is the ratio of a
point's density to the average density of its neighbours; scores
significantly greater than~1 indicate anomalies whose local neighbourhood
is denser than the point itself.

We use LOF in \emph{novelty} mode, fitting only on legitimate training data
(sub-sampled to 40{,}000 rows) and scoring new data.
$k \in \{10, 20, 50\}$ is swept with validation PR-AUC selection.

% ─── Reference Baselines ─────────────────────────────────────────────────
\subsubsection{Reference Baselines}
\label{sec:ref_baselines}

Two reference baselines bracket performance:
\begin{enumerate}
    \item \textbf{Random scores}: uniform $U(0,1)$ scores; PR-AUC should
          equal the base rate ($\approx 0.1\%$), serving as a sanity check.
    \item \textbf{$|\text{Amount}|$ score}: the absolute (scaled) transaction
          amount, testing whether high-value transactions are disproportionately
          fraudulent.
\end{enumerate}

% ─── Comparison to GMM ───────────────────────────────────────────────────
\subsubsection{Comparison to GMM}
\label{sec:compare_gmm}

The GMM (Sec.~\ref{sec:gmm}) is an explicit parametric density model:
it assumes the legitimate manifold follows a mixture of Gaussians and uses
the negative log-likelihood as the anomaly score.  This provides a
probabilistically principled score with well-understood calibration.

In contrast:
\begin{itemize}
    \item \textbf{Isolation Forest} is a non-parametric geometric method
          based on random axis-parallel partitions; it is robust to scaling
          and feature interactions but cannot estimate densities.
    \item \textbf{PCA reconstruction error} is a linear method that captures
          variance in the principal subspace; it fails when anomalies lie
          \emph{within} the principal subspace but differ in the residual.
    \item \textbf{OC-SVM} operates in kernel feature space and can model
          nonlinear boundaries, but is sensitive to kernel hyperparameters
          and scales poorly ($O(n^2)$--$O(n^3)$).
    \item \textbf{LOF} captures local density variations, making it effective
          for clusters of varying density, but is sensitive to the choice
          of $k$ and degrades in high dimensions.
\end{itemize}

Each approach has complementary failure modes: GMM struggles with heavy tails
and high dimensionality; Isolation Forest is weakened by irrelevant features
(axis-parallel splits); PCA requires anomalies to project into the minor
subspace; OC-SVM and LOF suffer from the curse of dimensionality.  Comparing
all approaches reveals which inductive biases are most appropriate for the
credit-card fraud dataset.
